\documentclass[11pt]{article}
% Portable preamble: compiles on Overleaf / any TeX Live with no conference .sty.
\usepackage[margin=1in]{geometry}
\usepackage{amsmath,amssymb,amsthm}
\usepackage{booktabs}
\usepackage{graphicx}
\usepackage{hyperref}
\usepackage{xcolor}
\usepackage{natbib}

\newtheorem{theorem}{Theorem}
\newtheorem{proposition}{Proposition}
\newtheorem{corollary}{Corollary}
\newcommand{\carriers}{\mathrm{car}}
\newcommand{\supp}{\mathrm{supp}}

\title{\textbf{SERUM}: Content-Aware Agentic Containment of Malware\\
Under an Unobserved Payload}
\author{Amritha S.}
\date{Draft --- \today}

\begin{document}
\maketitle

\begin{abstract}
Malware propagates only across hosts exploitable by its payload, so the
effective contagion graph is a payload-specific subgraph of the network. We
study the \emph{defender} who does not observe the payload and must infer
\emph{which exploit is spreading}, online, from the vulnerability-consistency
structure of the infected set, then allocate a limited containment budget under
that belief. We show (i) a combinatorial identifiability theory: because the
propagation subgraph is fixed by \emph{observable} node attributes (the CVEs each
host carries), the exploit is identifiable from a saturating outbreak iff the
intersection of the infected hosts' vulnerability profiles is a singleton---a
condition a defender can check in advance from its asset inventory; (ii) a
content-aware agent that plans under a Bayesian belief over the exploit,
optionally warm-started by a CVSS/LLM threat-intel prior; and (iii) an evaluation
on real NVD/CVE data with software-monoculture network zones. Against the
strongest deployable structure-only immunization the content-aware agent reduces
infection by $18.6\%$ ($p=1.1\times10^{-5}$), beats an oracle that receives the
signature after a delay, and is Pareto-dominant on availability and containment
speed. We are explicit about scope: it does not beat an oracle ensemble on
synthetic data, and the gated-propagation model itself is not novel; the
contribution is the coupling of \emph{observable-attribute identifiability} with
\emph{online, budgeted, content-aware containment}, grounded in real data.
\end{abstract}

\section{Introduction}
A self-propagating worm exploiting a specific vulnerability (CVE) can infect a
neighbouring host only if that host runs the vulnerable software. The network a
worm actually traverses is therefore not the physical topology but the
\emph{vulnerable subgraph} induced by the carriers of its target CVE. A defender
that knows which exploit is loose can defend exactly that subgraph; a defender
that sees only topology wastes budget immunizing hubs the payload cannot use.

The realistic difficulty is that the defender does \emph{not} observe the
payload. It sees who is infected and its own asset inventory. Because
propagation is vulnerability-gated, every host infected \emph{by spread} must
carry the target CVE---so the infected set is itself evidence about the exploit.
SERUM turns this into an online Bayesian inference (a POMDP whose hidden state is
the exploit) and acts under the resulting belief.

\paragraph{Contributions.}
\begin{itemize}
\item A combinatorial \textbf{identifiability theory} for the unobserved exploit
(\S\ref{sec:theory}): identifiable iff the intersection of infected profiles is a
singleton; confusability is exactly the subset order on carrier sets. Validated
against the Bayesian belief at $100\%$ agreement.
\item A \textbf{content-aware agent} (\S\ref{sec:method}) that plans under the
belief with a payload-conditioned exposed-vulnerable-degree score, switching from
isolation to precise patching as the belief sharpens, and an optional
CVSS/LLM-grounded threat-intel prior for cold starts.
\item A \textbf{real-data evaluation} (\S\ref{sec:exp}) on NVD/CVE with
software-monoculture zones, with paired statistics and an honest account of where
content-awareness does and does not help.
\end{itemize}

\paragraph{What is \emph{not} novel.} The gated-propagation model is a special
case of multitype bond percolation \citep{allard2009multitype} and the premise of
software-diversity worm models; classical worm models already gate on a
vulnerable population \citep{chen2003aawp,staniford2002own}. Learning-based
immunization is established \citep{meirom2021rlgn,fan2020finder}. We claim none of
these. The closest system, CyGym \citep{cygym2025}, has vulnerability-gated spread
and budgeted defense but uses a \emph{static} prior over zero-days; we make the
\emph{online inference of the exploit from the spread} the centerpiece.

\section{Problem formulation}\label{sec:formulation}
A host graph $G=(V,E)$; each host $v$ carries a vulnerability profile
$X(v)\subseteq C$ over a CVE universe $C$, with $\carriers(c)=\{v: c\in X(v)\}$
and vulnerable subgraph $G[c]$. A worm has unknown target $c^\star$. Propagation
is gated: an infected $u$ infects a susceptible neighbour $w$ with probability
$\beta$ only if $c^\star\in X(w)$ (SI; SIR with recovery time $r$ is a variant).
Each step the defender observes the infected set (not $c^\star$) and spends a
budget $B$ on \textsc{patch} (immunize a host, availability-preserving),
\textsc{isolate} (remove a host, availability cost), or \textsc{segment} (cut a
link), minimising a multi-objective cost of final infection, availability loss,
and time-to-containment.

\section{Identifiability of the exploit}\label{sec:theory}
Let $\supp(I)=\bigcap_{v\in I}X(v)$.

\begin{proposition}[Exact support]\label{prop:supp}
After observing a propagation-infected set $I$, the CVEs consistent with the
observation are exactly $\supp(I)$.
\end{proposition}

\begin{proposition}[Saturation]\label{prop:sat}
Let $R$ be the component of $G[c^\star]$ containing the seeds. As $\beta\to1$ the
outbreak infects exactly $R$, and the observer learns $\supp(R)$.
\end{proposition}

\begin{theorem}[Identifiability]\label{thm:id}
$c^\star$ is identifiable from a saturating outbreak over $R$ iff
$\bigcap_{v\in R}X(v)=\{c^\star\}$.
\end{theorem}

\begin{proposition}[Confusability = subset order]\label{prop:conf}
$c'$ remains consistent after a saturating $c$-outbreak iff $R\subseteq
\carriers(c')$; globally, $\carriers(c)\subseteq\carriers(c')$ implies $c'$ is
confusable with $c$ for every $c$-outbreak.
\end{proposition}

\begin{corollary}
Order CVEs by the subset relation on carrier sets. A CVE is globally identifiable
iff no other CVE's carrier set is a superset of its own.
\end{corollary}

Proofs are elementary set arguments (see supplementary \texttt{docs/THEORY.md}).
The consequence differs from cascade-mixture identifiability
\citep{hoffmann2020mixtures}: the candidate propagation graphs are not latent but
$\{G[c]\}$, fixed by observable profiles, so identifiability is a set-containment
condition a defender can \emph{evaluate in advance}. Empirically the Bayesian
belief converges to the truth exactly when Theorem~\ref{thm:id} predicts it
($116/116=100\%$); on real NVD profiles about $54\%$ of CVEs are identifiable
from a saturating outbreak.

\section{Method}\label{sec:method}
\textbf{Belief.} A posterior over $c^\star$: hard consistency (Prop.~1) for the
theory, and a \emph{soft} likelihood for robustness---each infected host
down-weights rather than eliminates a CVE, so a single false-positive detection
cannot exclude the truth. The prior is prevalence, or a CVSS/LLM
\emph{threat-intel} prior (weaponizability from severity $\times$ exploitability)
that warm-starts cold decisions.

\textbf{Content-aware policy.} Score each frontier host by its belief-weighted
\emph{exposed-vulnerable degree}---the expected number of susceptible,
exploitable neighbours it would infect next---and spend the budget top-down. When
the belief's support collapses, switch from isolation to patching (keep hosts
online). The full-observation oracle and an \emph{oracle-after-delay} (signature
arrives at $t=\delta$) bound the value of inference.

\section{Experiments}\label{sec:exp}
We build networks with segment-correlated software (connected zones share a
software image; a \texttt{homophily} knob controls monoculture) and draw target
CVEs from a prevalence band per trial. Every policy faces the identical outbreak
(paired design). Real-data profiles come from cleaned NVD records (network attack
vector, no user interaction); per-CVE $\beta$ from CVSS.

\begin{table}[t]\centering
\caption{Real NVD data, 60 paired outbreaks. Lower infection / higher
availability / faster containment is better.}
\begin{tabular}{lrrr}
\toprule
Policy & Infected & Availability & Contain@step \\
\midrule
no-defense & 14.0\% & 100.0\% & 10.5 \\
degree (best structural) & 1.5\% & 96.9\% & 3.1 \\
oracle-after-delay ($t{=}5$) & 1.4\% & 97.3\% & 2.9 \\
\textbf{content-aware (ours)} & \textbf{0.9\%} & \textbf{98.3\%} & \textbf{2.0} \\
content-aware oracle (bound) & 0.8\% & 100.0\% & 1.7 \\
\bottomrule
\end{tabular}
\end{table}

\paragraph{Findings.} Against the best fixed baseline, content-aware reduces
infection by $18.6\%$ (paired Wilcoxon $p=1.1\times10^{-5}$) and beats the
per-trial ensemble oracle on real data ($p=2\times10^{-3}$)---because real
software zones diverge from topological centrality. It beats oracle-after-delay:
inferring the exploit online beats waiting for a signature. The advantage scales
with outbreak severity and holds under SIR. The threat-intel prior gives a
$\sim37\%$ cold-start reduction under tight budgets and washes out with evidence.

\paragraph{Honest scope.} On \emph{synthetic} data content-aware does not beat an
oracle ensemble that cherry-picks the best structural heuristic per outbreak
($p\approx0.84$); and on near-universal (``popular'') exploits over
geometric/small-world topologies structure-only can win, since almost every host
is vulnerable. Content-awareness helps when the exploit is stealthy enough that
its victims are a distinctive, mislocated-from-the-hubs zone.

\section{Related work}
Source detection \citep{shah2011rumors} infers \emph{where} an outbreak started;
NetRate \citep{gomez2011netrate} infers edges/rates; SERUM infers \emph{which
exploit}. The nearest inference-from-shape is SCENARIOID \citep{scenarioid2023}
(offline scenario classification); we do online belief tracking tied to
identifiability. The POMDP/active-hypothesis-testing machinery is standard
\citep{naghshvar2013active,miehling2018pomdp}; our observation model
(vulnerability-gated cascade shape) and its coupling to identifiability are the
new part. Multi-objective immunization and active sensing
\citep{leskovec2007celf} and estimation-evasion games \citep{fanti2015spy}
inform the roadmap.

\section{Limitations and responsible use}
The propagation model is abstract (a probability over an exploitability graph);
inventories are assumed known; adversarial payload design against the inference
is future work. SERUM is defensive: it studies faster, lower-collateral
containment and contains no weaponizable attack code.

\bibliographystyle{plainnat}
\bibliography{refs}
\end{document}
